\section{Immutable Parent Pointers}
\label{sec:immutable-parent-pointers}

The immutable parent pointer is the foundational structural primitive of the Recursive Spawning Protocol. It is the single mechanism that transforms the delegation chain from a forensic reconstruction into a runtime-verifiable artifact, renders chain integrity a structural guarantee rather than a policy assertion, and makes all modes of autonomous drift architecturally impossible regardless of the operational urgency, architectural complexity, or adversarial sophistication of any machine actor in the chain.

This section specifies the parent pointer formally, defines the chain traversal operator using the recursive formulation, establishes the base case for the root human anchor, specifies the immutability enforcement mechanism through the Fact Layer write protocol, describes the Purpose Layer's spawn authorization role, and provides the complete chain traversal algorithm with correctness arguments.

% ----------------------------------------------------------------
\subsection{Formal Definition: ParentPointer}
\label{sec:parent-pointer-definition}

\begin{definition}[Parent Pointer]
\label{def:parent-pointer}
Let $\mathcal{N}$ be the set of all agent nodes in a \bxthree{} system implementing the Recursive Spawning Protocol. For any node $c \in \mathcal{N}$ spawned via a valid RSP spawn event, the relation $\ParentPointer{p}{c}$ is established exactly once at node provisioning via a Fact Layer atomic write, and satisfies four governing properties:

\begin{enumerate}
    \item[(P1) \textbf{Uniqueness}] For every child node $c \in \mathcal{N}$, there exists exactly one node $p \in \mathcal{N}$ such that $\ParentPointer{p}{c}$ holds at all times after provisioning. No child may have multiple simultaneous parent pointers. No mechanism exists by which an established pointer may be overwritten or redirected.

    \item[(P2) \textbf{Immutability}] After the provisioning write confirming $\ParentPointer{p}{c}$ is recorded in the Fact Layer ledger, no operation exists --- from any source, with any privilege level, under any system condition --- that can modify or delete $\ParentPointer{p}{c}$. The pointer is permanent for the operational lifetime of $c$. This is not a policy enforced by access control; it is a structural property of the protocol: the modification operation is not defined.

    \item[(P3) \textbf{Directionality}] The parent pointer is an intrinsic attribute of the child node, not a reference held by the parent node. The child $c$ carries $\ParentPointer{p}{c}$ as an immutable attribute in its own Fact Layer state record. The parent $p$ holds no mutable reference to $c$. A parent that has terminated, crashed, or been decommissioned does not affect $c$'s ability to reference its parent. The chain is traversable from any node outward to its complete lineage independent of the operational state of any other node.

    \item[(P4) \textbf{Human Anchor Termination}] The root of any RSP chain is a human principal node $h$ for which $\ParentPointer{\bot}{h}$ holds. The null parent marker $\bot$ indicates that $h$ is the ultimate delegation origin. No machine actor in the RSP model may serve as a chain root.
\end{enumerate}
\end{definition}

The distinction between P2-immutability and access-control-based immutability is architectural. An access-controlled pointer can be modified by any principal with sufficient privileges; immutability is a policy enforceable until credential theft, privilege escalation, or administrative compromise circumvents it. The Fact Layer write protocol makes immutability a structural property: there is no modification pathway defined, so no amount of privilege elevation can produce a valid write to an established $\ParentPointer{p}{c}$. The pointer is not protected --- it is structurally incapable of being changed.

\begin{prop}[Singleton Parent Property]
\label{prop:singleton-parent}
For any node $c \in \mathcal{N}$, there exists exactly one $p \in \mathcal{N}$ such that $\ParentPointer{p}{c}$ holds. No $c$ may simultaneously satisfy $\ParentPointer{p_1}{c}$ and $\ParentPointer{p_2}{c}$ for $p_1 \neq p_2$.
\end{prop}
\begin{pf}
Uniqueness follows directly from the parent pointer definition at provisioning: the Fact Layer atomic write records exactly one $(p, c)$ pair in the forensic ledger $\mathcal{L}$. The write protocol (Section~\ref{sec:fact-layer-write-protocol}) rejects any duplicate provisioning attempt for the same child $c$, and defines no operation for creating a second pointer. ∎
\end{pf}

% ----------------------------------------------------------------
\subsection{The Chain Operator}
\label{sec:chain-operator}

The chain operator $\Chain{\cdot}$ constructs the complete delegation lineage from any node to the human anchor by recursively applying the parent pointer function. This operator is the primary tool for chain verification, audit, and accountability attribution.

\begin{dfn}[Chain Operator]
\label{def:chain-operator}
For any agent node $a \in \mathcal{N}$ in a Recursive Spawning chain, the \emph{chain} $\Chain{a}$ is defined recursively as:

\begin{equation}
\Chain{a} =
\begin{cases}
\{\, \ParentPointer{\parent(a)}{a} \,\} \cup \Chain{\parent(a)} & \text{if } \parent(a) \neq \bot \\[6pt]
\{\, \ParentPointer{\bot}{a} \,\} & \text{if } \parent(a) = \bot \tag{base case}
\end{cases}
\end{equation}

\noindent where $\parent(a)$ denotes the unique parent of $a$ (i.e., the $p$ such that $\ParentPointer{p}{a}$ holds), and $\bot$ is the null parent marker designating the human anchor.
\end{dfn}

The recursive definition produces the complete ancestral chain: the immediate parent pointer of $a$, unioned with the chain of its parent, and so on, until the root human anchor is reached. The base case is the root node $h$ where $\parent(h) = \bot$, producing $\{\,\ParentPointer{\bot}{h}\,\}$.

\begin{corol}[Chain Termination]
\label{corol:chain-terminates}
For any provisioned node $a \in \mathcal{N}$, $\Chain{a}$ terminates in finite time at a human anchor node $h$ such that $\ParentPointer{\bot}{h}$ holds.
\end{corol}
\begin{pf}
The chain traversal follows the parent pointer relation $\alpha(a) = \parent(a)$ at each step. By the RSP depth bound (Section~\ref{sec:depth-bound}), the maximum chain depth is $D_{\max} + 1$, where $D_{\max}$ is the maximum number of permitted spawn layers. Each iteration moves strictly closer to the root (depth decreases by one), so the recursion cannot be infinite. The root is reached in at most $D_{\max} + 1$ steps, and by Property P4 the root is a human anchor. ∎
\end{pf}

\begin{dfn}[Chain Length]
\label{def:chain-length}
The \emph{length} of $\Chain{a}$, denoted $|\Chain{a}|$, is the number of nodes in the chain from $a$ (inclusive) to the human anchor $h$ (inclusive). By the base case definition, $|\Chain{h}| = 1$ for any human anchor $h$.
\end{dfn}

The chain length is at most $D_{\max} + 1$ by the depth bound, and at least $1$ for any node including the root human anchor. The chain is a total order under the parent-child relation, with a unique human anchor as the minimal element.

\begin{dfn}[Chain Depth]
\label{def:chain-depth}
The \emph{depth} of a node $a$, denoted $\mathrm{depth}(a)$, is $|\Chain{a}| - 1$: the number of parent-pointer steps from $a$ to the human anchor. Equivalently, $\mathrm{depth}(a) = k$ means $a$ is the $k$-th generation descendant of the human anchor.
\end{dfn}

By definition, $\mathrm{depth}(h) = 0$ for any human anchor $h$, and $\mathrm{depth}(c) = \mathrm{depth}(\parent(c)) + 1$ for any non-root node $c$.

% ----------------------------------------------------------------
\subsection{Immutability Enforcement: Fact Layer Write Protocol}
\label{sec:fact-layer-write-protocol}

The Fact Layer is the authoritative, append-only forensic ledger that records all RSP state transitions. It is the only mechanism capable of establishing or modifying the parent pointer relation $\ParentPointer{p}{c}$, and it enforces immutability through a write protocol that defines no modification pathway for an established pointer.

\medskip
\noindent\textbf{Fact Layer write protocol for parent pointer provisioning.}

When a Purpose Layer warrant $\phi$ authorizes a spawn event $\mathrm{spawn}(c, p, \sigma_c, t, \phi)$, the Fact Layer executes the following atomic write sequence:

\begin{enumerate}
    \item[(W1) \textbf{Record the parent pointer.} Write the tuple $(\ParentPointer{p}{c}, t, \mathrm{PROVISION})$ to the forensic ledger $\mathcal{L}$, where $t$ is the provisioning timestamp. This write is the sole authoritative record of the parent-child relation for $c$.

    \item[(W2) \textbf{Record the warrant reference.} Write the tuple $(\phi, \ParentPointer{p}{c}, t, \mathrm{AUTHORIZED})$ to $\mathcal{L}$, establishing the immutable link between the Purpose Layer warrant and the parent pointer it authorized. Both records are written in the same atomic transaction; neither exists without the other.

    \item[(W3) \textbf{Reject modification requests.} Define no code path for processing a modification request to an established $\ParentPointer{p}{c}$. Any operation attempting to overwrite, redirect, or delete $\ParentPointer{p}{c}$ is rejected at the protocol level before the Fact Layer attempts to execute it. There is no privilege level, no system condition, and no emergency override under which the Fact Layer will execute such a modification.

    \item[(W4) \textbf{Atomic coupling of pointer and warrant.} The parent pointer and its authorizing warrant are created in a single atomic transaction. There is no temporal window in which the pointer exists without authorization, or in which authorization exists without a corresponding pointer. This eliminates the class of exploits where a parent pointer is first established and then retroactively authorized.
\end{enumerate}

\medskip
\noindent\textbf{Why access control is insufficient.}

The distinction between access-control-based immutability and protocol-level immutability is the distinction between a promise and a guarantee. In an access-control model, immutability is enforced by restricting which principals may issue modification operations. If an attacker acquires sufficient privileges --- through credential theft, insider compromise, or software vulnerability --- the access control barrier is circumvented and immutability is violated. The immutability was a policy, enforceable until it was not.

In the Fact Layer write protocol, the immutability of $\ParentPointer{p}{c}$ is a property of the protocol itself. There is no operation defined for modifying $\ParentPointer{p}{c}$ after provisioning, so no amount of privilege elevation, credential compromise, or software vulnerability can produce a valid modification operation. The protocol does not enforce immutability --- it makes immutability the only valid state transition. The distinction is architectural, not a matter of degree.

This is the precise failure mode the Fact Layer write protocol eliminates: retroactive manipulation of chain topology for the purpose of accountability laundering. Without protocol-level immutability, an adversarial actor could re-parent a node to a favorable parent after the fact, obscuring the original authorization chain. With protocol-level immutability, this manipulation is structurally impossible.

\begin{prop}[Immutability of $\ParentPointer{p}{c}$]
\label{prop:parent-pointer-immutable}
For any child node $c$ with parent $p$, once the spawn event $\mathrm{spawn}(c, p, \sigma_c, t, \phi)$ has been recorded in the forensic ledger $\mathcal{L}$, the relation $\ParentPointer{p}{c}$ is fixed for the operational lifetime of $c$. No subsequent operation --- including operations issued by $p$, by any Purpose Layer actor, by the Bounds Engine, or by $c$ itself --- can alter $\ParentPointer{p}{c}$.
\end{prop}
\begin{pf}
The Fact Layer write protocol defines no code path for processing a modification request to an established $\ParentPointer{p}{c}$ (W1). All modification requests are rejected at the protocol level before the Fact Layer attempts to execute them (W3). Since there is no operation defined for modification, the predicate $\ParentPointer{p}{c}$ retains its original truth value indefinitely. ∎
\end{pf}

\begin{prop}[Atomic Pointer-Warrant Coupling]
\label{prop:atomic-pointer-warrant}
The parent pointer $\ParentPointer{p}{c}$ and its Purpose Layer warrant reference $\phi$ are created in a single atomic Fact Layer transaction. There exists no state in which $\ParentPointer{p}{c}$ is recorded without $\phi$, and no state in which $\phi$ is recorded without the corresponding $\ParentPointer{p}{c}$.
\end{prop}
\begin{pf}
Follows directly from W2 and W4: the atomic transaction condition in W4 requires both records to be written together; W2 explicitly links them in the same transaction. ∎
\end{pf}

% ----------------------------------------------------------------
\subsection{Spawn Authorization: Purpose Layer Role}
\label{sec:purpose-layer-authorization}

The Purpose Layer plays two structurally necessary roles in the parent pointer lifecycle: it originates the spawn authorization policy that defines the conditions under which a new parent pointer may be created, and it issues the spawn warrant that is the prerequisite for the Fact Layer provisioning write.

\medskip
\noindent\textbf{Warrant as prerequisite.}

The Fact Layer will not execute the provisioning write for $\ParentPointer{p}{c}$ unless a valid Purpose Layer warrant $\phi$ is presented at the time of the write. The warrant attests that the spawn event satisfies three authorization conditions:

\begin{enumerate}
    \item[(S1) \textbf{Human origin}: the chain root is a human principal $h$ (i.e., $\ParentPointer{\bot}{h}$ holds), and each successive parent is a correctly provisioned RSP node;

    \item[(S2) \textbf{Scope refinement}: the child node's operating scope $R(c)$ is a proper subset of the parent node's operating scope $R(p)$, ensuring delegation is a narrowing operation and not a lateral expansion;

    \item[(S3) \textbf{Bounds compliance}: the resulting chain depth does not exceed $D_{\max}$, ensuring the delegation chain remains within the maximum depth bound.
\end{enumerate}

These conditions are evaluated by the Purpose Layer at spawn time and attested to in the warrant $\phi$. The Fact Layer does not re-evaluate these conditions --- it trusts the Purpose Layer's attestation, sealed by the Purpose Layer's own cryptographic key. The immutability of the parent pointer that results from a valid warrant is guaranteed by the Fact Layer write protocol, not by the Purpose Layer.

\begin{dfn}[Spawn Warrant]
\label{def:spawn-warrant}
A spawn warrant $\phi$ for a spawn event $\mathrm{spawn}(c, p, \sigma_c, t, \phi)$ is a signed attestation from the Purpose Layer that conditions (S1), (S2), and (S3) hold for the proposed parent pointer $\ParentPointer{p}{c}$. The warrant is recorded atomically with the parent pointer in the Fact Layer forensic ledger (W2), forming an immutable authorization chain.
\end{dfn}

The Purpose Layer's role is authorization, not enforcement. Once the warrant is issued and recorded, the Fact Layer enforces immutability of the resulting pointer. The Purpose Layer cannot revoke a warrant after the Fact Layer records the pointer --- the atomic coupling property (Proposition~\ref{prop:atomic-pointer-warrant}) ensures both are created together or not at all, with no post-creation revocation pathway.

\medskip
\noindent\textbf{No self-bootstrapping.}

A machine actor cannot originate its own chain. Condition (S1) requires that every spawn event trace back through its parent chain to a human principal. A machine actor that attempts to spawn a child without a Purpose Layer warrant from a human-originated parent will fail at the Fact Layer write protocol: no warrant means no provisioning write, no parent pointer. The chain cannot be bootstrapped by a machine actor.

This is enforced structurally, not by policy. The Fact Layer's W1 condition requires a warrant; without one, the provisioning write does not occur. There is no fallback path, no emergency override, no administrative bypass. A machine actor without a Purpose Layer warrant from a human-originated chain cannot establish a parent pointer.

\begin{prop}[No Self-Originated Chains]
\label{prop:no-self-origination}
For any spawn event $\mathrm{spawn}(c, p, \sigma_c, t, \phi)$, if $\phi$ is a valid Purpose Layer warrant, then $p$ is either a human anchor or a correctly provisioned RSP node with a chain traceable to a human anchor. No machine actor can be the origin of an RSP chain.
\end{prop}
\begin{pf}
A valid warrant $\phi$ attests to (S1), which requires the parent chain to trace to a human principal. The Fact Layer requires $\phi$ as a prerequisite for provisioning (W1). Therefore, any provisioned parent pointer has a chain traceable to a human anchor, and no machine actor can be a chain origin. ∎
\end{pf}

% ----------------------------------------------------------------
\subsection{Chain Traversal Algorithm}
\label{sec:chain-traversal-algorithm}

The chain traversal algorithm is the runtime mechanism for computing $\Chain{a}$ and verifying its structural validity. Two algorithms are specified: Chain-Traverse, which computes the chain from any node to the human anchor, and Chain-Verify, which verifies that a chain satisfies all RSP structural constraints.

\bigskip
\noindent\textbf{Algorithm: Chain-Traverse.}

\begin{algorithm}
\caption{Chain-Traverse($a$) --- Compute the delegation chain from node $a$ to human anchor}
\label{alg:chain-traverse}
\begin{algorithmic}[1]
\REQUIRE $a$ is a provisioned RSP node
\ENSURE Ordered list of nodes from $a$ to the human anchor $h$

\STATE $chain \leftarrow []$
\STATE $current \leftarrow a$

\WHILE{$\alpha(current) \neq \bot$}
    \STATE $\text{append } current \text{ to } chain$
    \STATE $current \leftarrow \alpha(current)$
\ENDWHILE

\STATE $\text{append } current \text{ to } chain$ \COMMENT{Append the human anchor $\bot$-root}
\RETURN $chain$
\end{algorithmic}
\end{algorithm}

Chain-Traverse($a$) returns the ordered list $[a, \parent(a), \parent(\parent(a)), \ldots, h]$, where $h$ is the human anchor. The algorithm terminates because each iteration assigns $\alpha(current)$, which by the parent pointer immutability and the depth bound (Corollary~\ref{corol:chain-terminates}) reaches $\bot$ within at most $D_{\max} + 1$ iterations.

\begin{prop}[Chain-Traverse Correctness]
\label{prop:chain-traverse-correct}
For any provisioned node $a \in \mathcal{N}$, Chain-Traverse($a$) returns exactly $\Chain{a}$ as defined in Definition~\ref{def:chain-operator}.
\end{prop}
\begin{pf}
By induction on the number of iterations. Let $a_0 = a$, $a_{i+1} = \alpha(a_i)$. The loop invariant is that at the start of iteration $i$, $current = a_i$ and the chain contains $[a_0, a_1, \ldots, a_{i-1}]$. The loop appends $current$, then assigns $current \leftarrow \alpha(current) = a_{i+1}$. When $\alpha(current) = \bot$, the loop terminates after appending the human anchor. The resulting list is $[a_0, a_1, \ldots, a_k, h]$ where $a_k$ is the deepest node for which $\alpha(a_k) \neq \bot$ and $h$ satisfies $\alpha(h) = \bot$. By the chain operator definition, this is exactly $\Chain{a}$. ∎
\end{pf}

\bigskip
\noindent\textbf{Algorithm: Chain-Verify.}

Chain-Verify is the operational counterpart to the chain operator definition: it provides the runtime verification procedure that confirms a given node's chain satisfies all RSP structural constraints.

\begin{algorithm}
\caption{Chain-Verify($n$) --- Verify the authorization chain from node $n$ to human anchor}
\label{alg:chain-verify}
\begin{algorithmic}[1]
\REQUIRE $n$ is a provisioned RSP node
\ENSURE $\text{Boolean}$ indicating whether the chain is a structurally valid RSP chain

\STATE $chain \leftarrow \text{Chain-Traverse}(n)$
\STATE $m \leftarrow \text{len}(chain)$

\STATE \COMMENT{\textit{V1: Verify chain terminates at human anchor}}
\IF{$\alpha(chain[m-1]) \neq \bot$}
    \STATE \RETURN $\text{false}$ \COMMENT{Root is not $\bot$ --- malformed chain}
\ENDIF
\IF{$\neg h(chain[m-1])$}
    \STATE \RETURN $\text{false}$ \COMMENT{Root is not a designated human principal}
\ENDIF

\STATE \COMMENT{\textit{V2: Verify each spawn event has a Purpose Layer warrant}}
\FOR{$i \leftarrow 0$ \TO $m-2$}
    \STATE $child \leftarrow chain[i]$
    \STATE $parent \leftarrow chain[i+1]$
    \IF{$\neg \text{warrant\_exists}(child,\ parent)$}
        \STATE \RETURN $\text{false}$ \COMMENT{Missing Purpose Layer warrant for spawn}
    \ENDIF
\ENDFOR

\STATE \COMMENT{\textit{V3: Verify scope refinement at each link}}
\FOR{$i \leftarrow 0$ \TO $m-2$}
    \STATE $child \leftarrow chain[i]$
    \STATE $parent \leftarrow chain[i+1]$
    \IF{$\neg(R(child) \subset R(parent))$}
        \STATE \RETURN $\text{false}$ \COMMENT{Scope refinement constraint violated}
    \ENDIF
\ENDFOR

\STATE \COMMENT{\textit{V4: Verify depth bound compliance at each node}}
\FOR{$i \leftarrow 0$ \TO $m-1$}
    \IF{$\text{depth}(chain[i]) > D_{\max}$}
        \STATE \RETURN $\text{false}$ \COMMENT{Depth bound exceeded at node}
    \ENDIF
\ENDFOR

\RETURN $\text{true}$
\end{algorithmic}
\end{algorithm}

Chain-Verify returns \texttt{true} if and only if all four verification conditions hold simultaneously. V1 confirms the chain terminates at a human anchor (not at a machine actor or orphaned node). V2 confirms every spawn was authorized by the Purpose Layer (not self-bootstrapped by a machine actor). V3 confirms scope refinement was respected at every link (not lateral expansion beyond authorized intent). V4 confirms the depth bound was respected throughout (not approaching infinity).

The algorithm is deterministic and produces the same result regardless of which node in the chain initiates verification. There is no self-serving interpretation of authorization that a drifted or orphaned node could produce, because Chain-Verify inspects the complete chain topology and warrants in the Fact Layer ledger, not the node's own claims about its authorization. Chain verification is a property of the chain topology, not a property of the node being verified.

\begin{prop}[Chain-Verify Soundness]
\label{prop:chain-verify-soundness}
If Chain-Verify($n$) returns \texttt{true}, then $\Chain{n}$ is a valid RSP delegation chain terminating at a human principal, and all spawn events in the chain satisfy the Purpose Layer authorization conditions (S1), (S2), and (S3).
\end{prop}
\begin{pf}
The algorithm directly implements the RSP correctness conditions. V1 enforces P4 (human anchor termination). V2 enforces the Purpose Layer warrant prerequisite, which is equivalent to (S1). V3 enforces scope refinement, which is (S2). V4 enforces the depth bound, which is (S3). If all four conditions hold, the chain satisfies the RSP structural constraints by construction. ∎
\end{pf}

\begin{prop}[Chain-Verify Completeness]
\label{prop:chain-verify-completeness}
If $\Chain{n}$ is a valid RSP delegation chain terminating at a human principal, then Chain-Verify($n$) returns \texttt{true}.
\end{prop}
\begin{pf}
For a valid RSP delegation chain, each verification condition has a corresponding structural guarantee: the chain terminates at a human anchor (satisfying V1), each spawn event has a recorded Purpose Layer warrant (satisfying V2), scope refinement holds at each link by the spawn authorization condition (satisfying V3), and depth never exceeds $D_{\max}$ by the depth bound (satisfying V4). No condition triggers a \texttt{false} return, so the algorithm returns \texttt{true}. ∎
\end{pf}

\bigskip
\noindent\textbf{Computational cost.}

Chain-Traverse runs in $O(d)$ time where $d$ is the chain depth (at most $D_{\max} + 1$). Chain-Verify runs in $O(m)$ time for the traversal plus $O(m)$ for each of the three verification loops, yielding $O(m)$ total where $m$ is chain length (also bounded by $D_{\max} + 1$). The algorithms are linear in chain length and independent of the total number of nodes in the system.

The depth bound $D_{\max}$ provides a constant-time upper bound on chain traversal cost: the worst-case number of pointer dereferences required to reach the human anchor from any node is $D_{\max} + 1$, independent of the total number of nodes in the system. Mutable parent pointers would allow the chain to be extended retroactively, making the effective depth potentially unbounded and the worst-case escalation time unbounded as well. Immutability is what makes the depth bound a structural guarantee rather than a statistical one.

% ----------------------------------------------------------------
\subsection{Summary: Immutability as Architectural Constraint}
\label{sec:immutability-summary}

The immutable parent pointer is the load-bearing constraint of the Recursive Spawning Protocol. Without it, the chain is mutable --- subject to retroactive modification by actors with sufficient privilege --- and the RSP's correctness properties (drift prevention, chain integrity, termination at human anchor) collapse to policy conventions rather than structural guarantees.

The immutability is architectural, not a stronger access control policy, because it is enforced by the Fact Layer write protocol, which defines no valid modification operation for $\ParentPointer{p}{c}$ after provisioning. This is distinct from access-control-based immutability in four critical ways: there is no override path for any actor; the immutability holds under all system failure conditions including those that disable other enforcement mechanisms; all sources receive equivalent enforcement treatment regardless of privilege level; and the parent pointer and its Purpose Layer warrant are created atomically with no temporal window for inconsistency.

The chain operator $\Chain{a}$ is uniquely determined by the immutable parent pointers in the chain. The chain traversal algorithm computes this operator correctly and terminates. Chain-Verify provides the runtime verification procedure for chain structural validity. Together, these mechanisms make the delegation chain an immutable, verifiable, and auditable runtime artifact --- the foundation on which all other RSP correctness guarantees rest.